%!TEX root = ../mainquasiufficiale.tex
\chapter{Background}
\label{cap:background}

Questo capitolo introduce, da zero, ogni concetto usato nel resto della
tesi, così da poter leggere i Capitoli~\ref{cap:tsc-modelli} e
\ref{cap:confronto} senza assumere familiarità pregressa con nessuno dei
campi coinvolti. I quattro strumenti generali discussi (reti neurali,
apprendimento per rinforzo, reti neurali ricorrenti, e reti neurali a
grafo) saranno la base dei modelli proposti e analizzati in questo lavoro.

\section{Struttura del capitolo}
\label{sec:background-overview}

Il capitolo comprende cinque sezioni e questa sezione ne anticipa il contenuto.

Nella sezione \ref{sec:background-nn} si introducono i concetti minimi di rete
neurale (percettrone, strlayerati, funzione di loss, discesa del gradiente,
backpropagation) su cui ogni rete usata in questo studio, dal DQN alla
LSTM alle reti a grafo, è costruita.

Nella sezione \ref{sec:background-rl} si presenta il reinforcement learning,
dal paradigma agente-ambiente fino alle estensioni del DQN effettivamente
usate in questo lavoro, passando per la distinzione tra apprendimento
model-based e model-free e per l'apprendimento multi-agente. Il
reinforcement learning model-based, pur introdotto per completezza
definitoria, non è discusso oltre la sua definizione, in quanto nessun agente
in questa tesi apprende un modello esplicito della dinamica
dell'ambiente (§\ref{subsec:model-free}).

Nella sezione \ref{sec:background-rnn} si introduce la necessità di una memoria
per dati sequenziali e la cella LSTM che questo studio adotta per
integrare l'informazione temporale. Si discute brevemente anche la GRU,
un'alternativa più leggera, solo per contrasto, mentre in questa sede non viene utilizzata.

Nella sezione \ref{sec:background-gnn} si presenta il concetto di grafo
e il \emph{message passing} su cui ogni \emph{graph neural network} è
basata, con particolare attenzione alla GCN e alla GAT, le due
formulazioni confrontate sperimentalmente nel Capitolo~\ref{cap:confronto}
tramite le rispettive varianti meta-condizionate.

Infine, nella sezione \ref{sec:background-tsc} viene descritto il problema applicativo di questo lavoro:
il controllo semaforico, o traffic signal control. Si introduce il
vocabolario di base dello stesso, le tre generazioni di
metodi di controllo (a tempi fissi, reattivo, appreso) e si posiziona
esplicitamente MetaSTGAT, il modello di riferimento di questo lavoro,
all'interno di questa tassonomia.

\section{Reti Neurali}
\label{sec:background-nn}

Una \textbf{rete neurale} è una funzione parametrica che trasforma un input in un output attraverso una successione di funzioni matematiche, i cui parametri, chiamati pesi, vengono appresi dai dati invece di essere codificati manualmente.
L'unità base è il \emph{neurone artificiale} o \emph{percettrone}:
esso riceve un vettore di input $x = (x_1, x_2, \dots, x_n) \in \mathbb{R}^n$,
ne calcola una combinazione lineare pesata $z = w^\top x + b = \sum_{i=1}^{n} w_i x_i + b$
(con $w = (w_1, w_2, \dots, w_n) \in \mathbb{R}^n$ il vettore di pesi e $b \in \mathbb{R}$ un termine costante detto \emph{bias}),
e applica a $z$ una funzione non lineare $\sigma$ detta \emph{funzione di attivazione},
producendo l'output $\hat{y}=\sigma(z)$.

\begin{figure}[htbp]
  \centering
  \includegraphics[width=0.7\textwidth]{nn_percettrone.jpg}
  \caption{Il percettrone: gli input $x_1, x_2, x_3$ vengono combinati con i
    pesi $w_1, w_2, w_3$ e il bias $b$, la somma $z$ attraversa la funzione
    di attivazione, producendo l'output $\hat y$.}
  \label{fig:percettrone}
\end{figure}

Attivazioni comuni sono la sigmoide $\sigma(z) = 1/(1+e^{-z})$, che comprime $z$ in $(0,1)$,
la tangente iperbolica $\tanh(z) = (e^z - e^{-z})/(e^z + e^{-z})$, che lo comprime invece in $(-1,1)$,
e la ReLU (Rectified Linear Unit) $\max(0,z)$, che si limita ad azzerare i valori negativi.
La non linearità è essenziale in quanto senza di essa comporre più trasformazioni lineari come vedremo ora
produrrebbe comunque una singola trasformazione lineare, non riuscendo ad apprendere delle relazioni complesse e non lineari presenti nei dati.

\begin{figure}[htbp]
  \centering
  \includegraphics[width=0.85\textwidth]{nn_attivazioni.jpeg}
  \caption{Le tre funzioni di attivazione discusse: sigmoide, tangente
    iperbolica e ReLU.}
  \label{fig:attivazioni}
\end{figure}

Una rete neurale \emph{feedforward} o \emph{multi-layer perceptron} (MLP)
dispone i neuroni in \textbf{strati} (\emph{layer}) successivi: l'output
di ogni strato è l'input dello strato seguente, dall'\emph{input layer}
(i dati grezzi) attraverso uno o più \emph{hidden layer} fino
all'\emph{output layer} finale.
In forma matriciale, un singolo layer calcola
\[h^{(l+1)} = \sigma\!\left(W^{(l)} h^{(l)} + b^{(l)}\right)\]
dove
$h^{(l)}$ raccoglie gli output di tutti i neuroni dello strato $l$,
$W^{(l)}$ è la matrice che impila i pesi di ciascun neurone di quello strato e
$b^{(l)}$ è il vettore dei bias.
Impilare più strati permette alla rete di comporre trasformazioni via via più astratte dell'input originale.

\begin{figure}[htbp]
  \centering
  \includegraphics[width=0.6\textwidth]{nn_mlp.png}
  \caption{Una rete feedforward con un input layer a 4 nodi, un hidden
    layer a 5 nodi e un output layer a un nodo: ogni connessione porta con
    sé un peso appreso, omesso per chiarezza.}
  \label{fig:mlp}
\end{figure}

\paragraph{Addestramento: discesa del gradiente e backpropagation.}
Una rete neurale viene inizializzata con pesi e bias casuali e viene addestrata
iterativamente confrontando l'output prodotto con quello desiderato.
Una funzione di loss $\mathcal{L}(\hat y, y)$ quantifica la distanza tra l'output prodotto $\hat y$ e quello desiderato
$y$ (che sia esso una etichetta di classificazione o un valore reale nel caso della regressione).
L'addestramento consiste nel regolare ogni peso nella direzione che riduce $\mathcal{L}$, un'operazione
nota come \emph{discesa del gradiente}:
\[w \leftarrow w - \eta \frac{\partial \mathcal{L}}{\partial w},\]
dove $\eta$, il \emph{tasso di apprendimento} o \emph{learning rate}, controlla l'ampiezza di ogni aggiornamento.
Calcolare la derivata parziale $\partial \mathcal{L}/\partial w$ per ogni peso della rete, dal più vicino
all'output al più lontano, richiede applicare ripetutamente la \emph{chain rule} del calcolo differenziale attraverso tutti gli strati
intermedi: l'algoritmo che esegue questo calcolo in modo efficiente, una
sola passata dall'output verso l'input, è la \emph{backpropagation}.
Nella pratica, il gradiente non è esatto ma si stima su un piccolo sottoinsieme casuale di
dati (\emph{mini-batch}) alla volta invece che sull'intero insieme di
dati disponibile, una variante nota come discesa del gradiente
stocastica, molto più economica computazionalmente a fronte di una stima
del gradiente più rumorosa ma sufficiente a guidare comunque
l'addestramento verso una soluzione utile.

\paragraph{Oltre la discesa del gradiente pura.} La regola di
aggiornamento vista sopra usa lo stesso tasso di apprendimento $\eta$ per
ogni peso della rete, un limite quando la superficie di loss ha curvatura
molto diversa da un peso all'altro (alcune direzioni richiederebbero un
passo grande, altre uno minimo): un ottimizzatore più sofisticato adatta
l'ampiezza dell'aggiornamento parametro per parametro, sulla base della
storia recente dei gradienti osservati. Il \emph{momentum} accumula una
media mobile esponenziale dei gradienti passati e aggiorna i pesi in quella
direzione media invece che nella sola direzione del gradiente corrente,
smorzando le oscillazioni lungo le direzioni in cui il segno del gradiente
cambia di continuo. \textbf{RMSprop} \citep{tieleman2012rmsprop} adotta un
principio complementare: mantiene per ciascun peso una media mobile del
quadrato dei gradienti recenti e divide l'aggiornamento per la sua radice,
riducendo il passo sui pesi il cui gradiente è stato storicamente grande in
valore assoluto e aumentandolo su quelli il cui gradiente è stato piccolo,
così da rendere il progresso più uniforme tra pesi con scale molto diverse.
\textbf{Adam} \citep{kingma2015adam}, un altro ottimizzatore molto usato combina invece entrambe le idee, media mobile
del gradiente e del suo quadrato. Questo lavoro
adotta \textbf{RMSprop}, coerentemente con il modello di riferimento
\citep{wang2022metastgat}, per l'addestramento di ogni rete $Q$ descritta
nel Capitolo~\ref{cap:tsc-modelli}.

Ogni rete usata in questo lavoro, dal DQN (§\ref{sec:background-rl}) alla
LSTM (§\ref{sec:background-rnn}) alle reti a grafo
(§\ref{sec:background-gnn}), è in ultima analisi una particolare
composizione di questi stessi due ingredienti: strati di neuroni con
un'attivazione non lineare, addestrati per discesa del gradiente tramite
backpropagation. Ciò che cambia da un'architettura all'altra è come gli
strati sono collegati tra loro, non il principio di addestramento
sottostante.

\section{Reinforcement Learning}
\label{sec:background-rl}

\subsection{Il paradigma agente-ambiente}

L'apprendimento per rinforzo o \emph{reinforcement learning} (RL) è un
paradigma di apprendimento automatico in cui un \emph{agente} impara a
un certo comportamento interagendo ripetutamente con un \emph{ambiente}, senza che
gli siano forniti esempi di comportamento corretto. A ogni passo di
tempo $t$, l'agente osserva lo stato corrente dell'ambiente, sceglie
un'\emph{azione} tra quelle disponibili, e riceve in cambio due informazioni:
un nuovo stato (l'ambiente è cambiato, anche per effetto dell'azione
scelta) e un numero scalare, la \emph{ricompensa}, che indica quanto
quell'azione sia stata desiderabile in quel momento.
Il problema che il RL affronta è per l'appunto
questo, cioè imparare per tentativi ripetuti quale comportamento massimizzi
la somma delle ricompense ricevute nel tempo, non solo quella immediata.
La Figura~\ref{fig:rl-loop} riassume questo ciclo, utile come riferimento
per la terminologia introdotta nel resto della sezione.

\begin{figure}[htbp]
  \centering
  \includegraphics[width=0.55\textwidth]{rl_loop}
  \caption{Il ciclo agente-ambiente: la policy sceglie un'azione,
    l'ambiente risponde con un nuovo stato e una ricompensa, e
    l'algoritmo di apprendimento usa quella ricompensa per aggiornare la
    policy.}
  \label{fig:rl-loop}
\end{figure}

Vale la pena distinguere questo paradigma da quello dell'\emph{apprendimento
  supervisionato}. In quest'ultimo una rete neurale riceve coppie input-output correttamente associate e impara
a riprodurre quella corrispondenza. Nel RL non esiste alcuna etichetta
corretta da imitare: l'agente deve scoprire da solo, per tentativi, quale
azione sia buona in ciascuno stato, guidato solo dalla
ricompensa ricevuta in seguito alla azione che l'ha causata.
È un problema strutturalmente più difficile,
ma anche l'unico adatto a situazioni, come il controllo semaforico, in
cui nessuno conosce a priori quale sia la "decisione corretta" da
etichettare.

\subsection{Processo decisionale di Markov}

La formalizzazione matematica standard di un problema di RL è il
\emph{processo decisionale di Markov} (MDP) \citep{suttonbarto2018},
definito dalla tupla
\[\mathcal{M} = (\mathcal{S}, \mathcal{A}, P, R, \gamma)\]
dove $\mathcal{S}$ è l'insieme di tutti gli stati possibili dell'ambiente,
$\mathcal{A}$ l'insieme di tutte le azioni possibili, $P(s' \mid s,a)$ la
probabilità di transitare nello stato $s'$ avendo eseguito l'azione $a$
nello stato $s$, $R(s,a)$ la ricompensa attesa per quella stessa coppia, e
$\gamma \in [0,1)$ il \emph{fattore di sconto}, che pondera l'importanza
delle ricompense future rispetto a quelle immediate: un $\gamma$ vicino a
$0$ rende l'agente "miope", interessato quasi solo alla ricompensa
immediata, un $\gamma$ vicino a $1$ lo rende attento anche a conseguenze
lontane nel tempo. Questa formalizzazione si basa sull'assunzione di Markov, per la quale la
probabilità di transizione $P(s'\mid s,a)$ dipende solo dallo stato
corrente $s$ e non dalla sequenza di stati che ha portato allo steso.

Un agente segue una \emph{politica} $\pi(a \mid s)$, una distribuzione
di probabilità sulle azioni condizionata allo stato corrente; il suo
obiettivo è massimizzare il \emph{ritorno} atteso
\[G_t = \sum_{k=0}^{\infty} \gamma^k r_{t+k}\]
ovvero la somma scontata delle ricompense future a partire dall'istante
$t$. Per poter ragionare su quanto buona sia una politica prima ancora di
eseguirla fino in fondo, si definiscono due funzioni valore: la
funzione valore-stato $V^\pi(s) = \mathbb{E}_\pi[G_t \mid s_t = s]$ e la
funzione valore-azione $Q^\pi(s,a) = \mathbb{E}_\pi[G_t \mid s_t=s,
    a_t=a]$, che misurano rispettivamente quanto ritorno atteso si ottiene
seguendo $\pi$ da uno stato, o da una coppia stato-azione. Queste
soddisfano l'\emph{equazione di Bellman}, che esprime il valore di uno
stato in funzione del valore atteso degli stati successivi:
\begin{equation}
  V^\pi(s) = \sum_a \pi(a\mid s) \sum_{s'} P(s'\mid s,a) \big[R(s,a) + \gamma V^\pi(s')\big]
\end{equation}
Per una politica ottima $\pi^*$ che massimizza $V^\pi(s)$ per ogni stato
simultaneamente, vale l'analoga equazione di Bellman di ottimalità:
\begin{equation}
  Q^*(s,a) = \mathbb{E}_{s' \sim P(\cdot \mid s,a)}\left[
    R(s,a) + \gamma \max_{a'} Q^*(s', a')
    \right]
  \label{eq:bellman-optimal}
\end{equation}
Se si conoscesse $Q^*$, la politica ottima sarebbe immediata da ricavare:
scegliere sempre $\operatorname*{argmax}_a Q^*(s,a)$. Il problema del RL
si riduce quindi, in buona parte, a stimare $Q^*$ senza conoscerlo a
priori.

\subsection{Apprendimento model-based e model-free}
\label{subsec:model-free}

Esiste una distinzione fondamentale tra due famiglie di algoritmi RL, a
seconda che l'agente costruisca o meno una rappresentazione esplicita
della dinamica dell'ambiente. Un algoritmo \emph{model-based} apprende
o riceve un modello approssimato di $P(s'\mid s,a)$ e $R(s,a)$, e lo usa
per simulare le conseguenze di più
sequenze di azioni prima di sceglierne una, senza eseguirle davvero nell'ambiente.
Un algoritmo \emph{model-free}, al contrario, non costruisce alcun modello della dinamica dell'ambiente, bensì stima
direttamente $Q^*$ dall'esperienza osservata, senza mai
provare a prevedere esplicitamente lo stato successivo. Il vantaggio del
model-free è la semplicità e la minore dipendenza dalla qualità di un
modello che potrebbe essere difficile da apprendere accuratamente; lo
svantaggio è un uso tipicamente meno efficiente dell'esperienza raccolta,
poiché non la si riusa per simulare scenari mai osservati direttamente.

\textbf{Questo lavoro adotta un approccio interamente model-free}. Né il
modello di riferimento \citep{wang2022metastgat} né le varianti
implementate in questa tesi (Capitolo~\ref{cap:tsc-modelli}) apprendono o
usano un modello esplicito della dinamica del traffico. L'agente stima
direttamente una funzione $Q$ dall'esperienza simulata in CityFlow, con una
famiglia di tecniche descritte nel resto di questa sezione.

\subsection{Apprendimento online e offline}
\label{subsec:online-offline}

Una distinzione ortogonale a quella model-based/model-free appena vista
riguarda da dove arriva l'esperienza usata per apprendere.

Un algoritmo \emph{online} interagisce ripetutamente con l'ambiente
durante l'addestramento, raccogliendo nuova esperienza con la politica
corrente e usandola subito per migliorarsi, in un ciclo continuo di raccolta e apprendimento.
Un algoritmo \emph{offline}, al contrario, apprende da un
insieme fisso di transizioni raccolto in precedenza, magari con una
politica diversa o addirittura sconosciuta, senza alcuna ulteriore
interazione con l'ambiente durante il training. Questo risulta essere un approccio necessario
quando interagire con l'ambiente reale è costoso o rischioso.
Un dettaglio che vale la pena chiarire subito è che un
\emph{replay buffer} come quello che il DQN (descritto a breve) introduce,
non lo rende di per sé un metodo offline. Il buffer infatti si riempie
continuamente di transizioni generate via via da una politica in
evoluzione, con raccolta dati e apprendimento che procedono in parallelo, non in
due fasi separate come negli algoritmi offline.

\textbf{Questo lavoro è interamente online}, ogni transizione
usata per un aggiornamento proviene dalla simulazione in corso in
\textsc{CityFlow}, mai da un insieme di dati raccolto in anticipo.

\subsection{Ampiezza e profondità dell'aggiornamento}
\label{subsec:backup-space}

Una seconda distinzione, anch'essa ortogonale alle precedenti, riguarda il
modo in cui si aggiorna una stima di valore combinando una ricompensa con
la stima del valore in uno o più stati successivi. Due assi indipendenti
classificano ogni aggiornamento, come mostrato in
Figura~\ref{fig:rl-backup-space}.

Il primo asse riguarda l'\textbf{ampiezza} dell'aggiornamento. Un
aggiornamento completo (\emph{full backup} in figura) calcola
l'aspettativa matematica considerando tutte le azioni possibili e tutti i
possibili stati successivi $s'$, pesati dalla probabilità di transizione
$P(s'\mid s,a)$ fornita da un modello noto dell'ambiente.
Un aggiornamento campionato (\emph{sample backup} in
figura) si basa invece su una singola traiettoria o un'esperienza realmente osservata,
senza calcolare tutte le diramazioni possibili.

Il secondo asse riguarda la \textbf{profondità} dell'aggiornamento. Un
aggiornamento superficiale, o a corto raggio (\emph{shallow backup} in
figura), guarda solo a pochissimi passi nel futuro, nel caso di TD(0) un
solo passo d'azione, e sostituisce il ritorno futuro con
la stima già disponibile dell'ultimo passo considerato. Questa è una tecnica nota come
\emph{bootstrapping}, per cui si costruisce una stima a partire da
un'altra stima invece che dall'esito osservato per intero. Un
aggiornamento profondo, o a lungo raggio (\emph{deep backup} in figura),
tiene conto delle conseguenze fino alla fine
dell'episodio, usando solo ricompense effettivamente realizzate e non stimate.

\begin{figure}[htbp]
  \centering
  \includegraphics[width=0.68\textwidth]{rl_backup_space}
  \caption{Lo spazio dei metodi di apprendimento e pianificazione,
    classificati secondo la profondità dell'aggiornamento (asse
    orizzontale, superficiale a sinistra, profondo a destra) e la sua
    ampiezza (asse verticale, completo in alto, campionato in basso); le
    etichette in figura usano i termini inglesi \emph{shallow/deep
      backups} per l'asse orizzontale e \emph{full/sample backups} per il
    verticale. I quattro angoli corrispondono a dynamic programming,
    ricerca esaustiva, Temporal Difference e Monte Carlo. Adattata da
    \citet{suttonbarto2018}.}
  \label{fig:rl-backup-space}
\end{figure}

Le combinazioni dei due assi individuano quattro famiglie di metodi:

\begin{itemize}
  \item \textbf{Completo e superficiale}: la \emph{dynamic programming},
        che calcola l'aspettativa esatta su tutti i successori $s'$ ma si
        ferma dopo un solo passo, facendo bootstrapping sulla stima già
        disponibile del valore di $s'$ invece di espandere ulteriormente
        l'albero.
  \item \textbf{Completo e profondo}: la \emph{ricerca esaustiva}, che
        come la dynamic programming pesa ogni successore con la sua
        probabilità, ma che espande però l'intero
        albero delle conseguenze fino alla fine dell'episodio o a un
        orizzonte fissato. È il caso ideale in cui ci si potesse
        permettere di enumerare ogni possibile continuazione.
  \item \textbf{Campionato e profondo}: il \emph{Monte Carlo} non
        usa $P(s'\mid s,a)$, osserva un'unica traiettoria
        effettivamente realizzata e ne percorre per intero le conseguenze
        fino alla fine dell'episodio, calcolando il ritorno effettivo
        $G_t = \sum_k \gamma^k r_{t+k}$ senza mai usare bootstrapping.
  \item \textbf{Campionato e superficiale}: il \emph{Temporal
          Difference} (TD), che come Monte Carlo osserva una singola
        transizione campionata anziché un'aspettativa sul modello, ma
        come la dynamic programming si ferma dopo un solo passo,
        facendo bootstrapping sulla stima corrente del valore dello stato
        successivo.
\end{itemize}

Nella sua versione più semplice di Temporal Difference, TD(0),
l'aggiornamento della funzione valore-stato è
\[V(s_t) \leftarrow V(s_t) + \eta\big[(r_t + \gamma V(s_{t+1})) - V(s_t)\big],\]
dove il termine tra parentesi quadre è il \emph{Temporal Difference
  error}, ovvero la discrepanza tra la stima corrente e un target
calcolato guardando avanti di un solo passo. Il bootstrapping ha un costo,
l'errore della stima usata come target si propaga nell'aggiornamento, ma
un beneficio pratico decisivo rispetto a Monte Carlo: l'agente può
aggiornarsi a ogni singolo passo, senza attendere la fine dell'episodio,
un requisito essenziale per imparare da compiti molto lunghi o
addirittura privi di una fine naturale. Il Q-learning descritto di
seguito applica esattamente questo principio alla funzione valore-azione
$Q$ invece che alla sola funzione valore-stato $V$.

\textbf{Questo lavoro si colloca nella casistica dell'aggiornamento campionato e superficiale, quello del Temporal Difference}.
L'assenza di un modello della dinamica del traffico (§\ref{subsec:model-free}) esclude
già di per sé la dynamic programming e la ricerca esaustiva, entrambe
legate a un aggiornamento completo; inoltre la simulazione in \textsc{CityFlow} procede
per episodi lunghi e continui su intersezioni multiple, un contesto in
cui attendere la fine dell'episodio per ogni aggiornamento, come farebbe
Monte Carlo, sarebbe sia lento sia inutilmente rumoroso rispetto al
bootstrap a un passo. Il Q-learning e il DQN che ne deriva, discussi di
seguito, adottano infatti un aggiornamento TD, campionato e superficiale.

\subsection{Q-learning e Deep Q-Network}

\paragraph{Q-learning tabellare.} Il \textbf{Q-learning} approssima $Q^*$
applicando iterativamente l'Equazione~\ref{eq:bellman-optimal} come
regola di aggiornamento tramite bootstrapping: a ogni transizione
osservata $(s,a,r,s')$, aggiusta il valore stimato $Q(s,a)$ verso il
target $r + \gamma \max_{a'} Q(s',a')$, la stima corrente al passo
successivo, invece di attendere il ritorno realizzato per intero a fine
episodio come farebbe un metodo Monte Carlo. Il nome "tabellare" indica
che questi aggiornamenti possono essere eseguiti considerando una tabella con una riga per ogni stato e una
colonna per ogni azione. L'aggiornamento modifica una voce alla volta e
ripetendo l'operazione su molte transizioni, ogni voce $Q(s,a)$ converge verso il
vero ritorno atteso dell'eseguire $a$ in $s$. Si ottiene in questo modo unaa politica
ottima, senza il bisogno di indicazioni esplicite sulle azioni da intraprendere.

L'approccio tabellare ha però un limite rappresentato dalle dimensioni della tabella.
Se lo spazio degli stati e lo spazio delle azioni sono continui o ad alta dimensionalità il numero di combinazioni possibili esplode,
rendendo questo approccio impraticabile.
Questo è il caso del controllo semaforico multi-agente, dove la dimensione dello spazio degli stati rende una
tabella esplicita irrealizzabile sia da costruire sia da aggiornare.

\paragraph{Deep Q-Network (DQN).} \citet{mnih2015dqn} risolvono il
problema sostituendo la tabella con una rete neurale $Q(s,a;\theta)$
allenata ad approssimarne i valori: invece di una voce indipendente per
ogni coppia $(s,a)$, la rete apprende una funzione continua che
generalizza a stati mai visti interpolando tra stati simili osservati
durante l'addestramento. I pesi $\theta$ vengono aggiornati
per discesa del gradiente in modo che $Q(s,a;\theta)$ si avvicini allo
stesso target di Bellman usato dal Q-learning tabellare, con la
differenza che l'aggiornamento modifica l'intera rete invece di una
singola voce, propagando l'informazione anche a stati simili mai
osservati direttamente. È esattamente questa proprietà a rendere il DQN
necessario in questo lavoro: lo spazio di stato del controllo semaforico
multi-agente è troppo grande per una tabella, ma una rete neurale può
generalizzare a partire da un numero gestibile di transizioni osservate
durante la simulazione. Per mitigare l'instabilità
dell'addestramento dovuta alla correlazione temporale dei campioni (stati
consecutivi di uno stesso episodio sono molto simili tra loro, violando
l'assunzione di campioni indipendenti richiesta dalla discesa del
gradiente stocastico), gli autori introducono inoltre il \emph{replay
  buffer}: una memoria in cui si accumulano le transizioni $(s,a,r,s')$
osservate e da cui si campionano poi mini-batch in ordine casuale per
l'addestramento, rompendo la correlazione temporale. Come anticipato
(§\ref{subsec:online-offline}), il buffer non cambia la natura online del
metodo, in quanto si riempie di transizioni raccolte dalla politica corrente mentre
l'addestramento procede, non da un insieme fisso raccolto in anticipo.

\subsection{Esplorazione $\varepsilon$-greedy}

Un agente completamente \emph{greedy}, che segue cioè sempre e solo l'azione a valore $Q$ stimato più
alto, non raccoglierebbe mai esperienza sulle conseguenze di azioni
alternative, rischiando di convergere su una politica sub-ottima per via
di una mancata esplorazione. Ad esempio, se nelle prime fasi
dell'addestramento un'azione appare leggermente migliore di un'altra, un
agente completamente \emph{greedy} la prediligerebbe sistematicamente,
impedendo di fatto all'agente di scoprire il potenziale di strategie alternative.
La tecnica più comune per bilanciare sfruttamento (\emph{exploitation}, scegliere
l'azione ritenuta migliore) ed esplorazione (\emph{exploration},
raccogliere informazione su azioni meno note) è $\varepsilon$-greedy: a
ogni passo di decisione, con probabilità $\varepsilon$ l'agente sceglie
un'azione uniformemente a caso, e con probabilità $1-\varepsilon$ sceglie
$\operatorname*{argmax}_a Q(s,a;\theta)$. Il valore di $\varepsilon$
tipicamente decresce nel corso dell'addestramento, da un valore iniziale
vicino a 1 (quasi solo esplorazione, quando la stima di $Q$ è ancora
inaffidabile) verso un valore finale piccolo ma non nullo (quasi solo
sfruttamento, mantenendo comunque un minimo di esplorazione residua).

\subsection{Estensioni usate in questo lavoro}

Il DQN di base ammette diverse estensioni indipendenti, ciascuna delle
quali corregge una specifica fonte di instabilità o inefficienza; la
procedura di training descritta nel Capitolo~\ref{cap:tsc-modelli} le
adotta tutte, motivandole caso per caso rispetto all'articolo di
riferimento:
\begin{itemize}
  \item Il \textbf{Double DQN} \citep{vanhasselt2016double} disaccoppia la
        selezione dell'azione dalla sua valutazione, usando la rete online per
        scegliere l'azione migliore e una \emph{target network}
        $Q(\cdot;\theta^-)$, una copia periodicamente sincronizzata della rete
        online, solo per valutarla:
        $y = r + \gamma\, Q\!\left(s', \operatorname*{argmax}_{a'}
          Q(s',a';\theta);\ \theta^-\right)$. Questo meccanismo evita il problema
        del \emph{moving target}, che si verifica ottimizzando una stima
        rispetto a parametri in continua variazione, e riduce il bias
        sistematico di sovrastima del DQN standard che si verifica quando si usa la stessa rete sia
        per scegliere sia per valutare l'azione.
  \item Il \textbf{Prioritized Experience Replay} (PER,
        \citealp{schaul2016per}) sostituisce il campionamento uniforme dal
        replay buffer con un campionamento proporzionale al \emph{Temporal Difference error} $\delta_i$.
        Le transizioni che la rete predice peggio vengono ricampionate più
        spesso, concentrando la capacità di apprendimento dove serve di più.
        Assegnando a ciascuna transizione una priorità $p_i = |\delta_i| +
          \varepsilon$, la probabilità di campionamento è $P(i) = p_i^\alpha /
          \sum_k p_k^\alpha$, con $\alpha \in [0,1]$ iperparametro che ne
        controlla la dipendenza dall'errore. Per correggere il bias introdotto
        da un campionamento non uniforme, la loss viene pesata tramite
        coefficienti di \emph{importance sampling} $w_i = \left(N \cdot
          P(i)\right)^{-\beta}$, dove $N$ è la dimensione del buffer e $\beta \in
          [0,1]$ iperparametro dinamico che ne controlla il grado di correzione.
  \item La \textbf{Huber loss} (o \emph{smooth L1}) sostituisce nel caso di valori continui la funzione di loss data dall'errore
        quadratico medio (MSE) con una funzione quadratica vicino allo zero e
        lineare lontano da esso, risultando meno sensibile agli outlier
        prodotti da transizioni con errore di stima molto grande (con MSE, un
        singolo errore enorme domina il gradiente dell'intero batch; con Huber
        il suo contributo cresce solo linearmente):
        \begin{equation}
          L_\delta(y, \hat y) =
          \begin{cases}
            \tfrac{1}{2}(y-\hat y)^2                           & \text{se } |y-\hat y| \le \delta \\
            \delta\left(|y-\hat y| - \tfrac{1}{2}\delta\right) & \text{altrimenti}
          \end{cases}
        \end{equation}
        Tipicamente $\delta = 1$.
  \item Quando la rete $Q$ include un modulo ricorrente
        (§\ref{sec:background-rnn}), l'assunzione di Markov sulle singole
        transizioni campionate dal replay buffer viene invalidata: lo stato
        ricorrente $(h,c)$ usato durante il training viene inizializzato a
        zero, introducendo un disallineamento rispetto allo stato generato
        sequenzialmente durante la raccolta dati, un problema noto come
        \emph{hidden state staleness}. Seguendo \citet{kapturowski2019r2d2}, il campionamento si esegue su
        intere sequenze di lunghezza $L$, dedicando i primi passi
        della sequenza (fase di \emph{burn-in}) esclusivamente alla ricostruzione di uno stato
        ricorrente plausibile, senza calcolare gradiente, per poi eseguire la
        \emph{backpropagation through time} (BPTT, §\ref{sec:background-rnn})
        sui passi restanti.
\end{itemize}

\subsection{Apprendimento multi-agente}
\label{subsec:marl}

Quanto descritto finora assume un solo agente che interagisce con
un solo ambiente. Quando più agenti operano nello stesso ambiente,
ciascuno osservando solo una porzione dello stato globale e influenzando
con le proprie azioni anche ciò che gli altri agenti osservano e
ricevono come ricompensa, il problema si generalizza da MDP a
\emph{processo decisionale di Markov decentralizzato e parzialmente
  osservabile} (Dec-POMDP). Il termine "decentralizzato" si riferisce al fatto che non esiste un'unica
azione scelta da un solo decisore, ma $N$ azioni scelte
indipendentemente da $N$ agenti; "parzialmente osservabile" indica invece che ogni
agente vede solo il proprio stato locale, non lo stato globale
completo dell'ambiente. Nel contesto di questo lavoro, la decentralizzazione e la parzializzazione
si concretizzano nel campo visivo limitato di ogni intersezione (Capitolo~\ref{cap:tsc-modelli}), che esclude deliberatamente
dallo stato osservato ciò che accade agli altri incroci.

Il multi-agente introduce una difficoltà che il caso a singolo agente non
ha: la \emph{non stazionarietà}. L'equazione di Bellman e l'intera teoria
del Q-learning assumono che l'ambiente risponda in modo coerente nel
tempo alla stessa azione nello stesso stato; ma se più agenti apprendono
simultaneamente, ciascuno sta anche cambiando il proprio comportamento
mentre gli altri lo osservano, cosicché per ogni singolo agente
l'ambiente (che include il comportamento degli altri agenti) appare
cambiare regole nel tempo, anche se le vere dinamiche fisiche sottostanti
restano identiche.

Una seconda difficoltà è
l'\emph{assegnazione del merito} (\emph{credit assignment}), per la quale quando più
agenti contribuiscono a un solo risultato osservato, non è ovvio quale
frazione di quel risultato attribuire a ciascuno. Tale problema però non si pone se ogni agente riceve una ricompensa
calcolata esclusivamente sul proprio stato locale, così come avviene in questo lavoro.
Una versione attenuata resta comunque presente: lo stato locale di un agente, e
quindi la sua stessa ricompensa, dipende in parte dalle scelte degli altri agenti, per cui
l'effetto della propria azione e quello dell'azione altrui rimangono in
parte intrecciati anche con una ricompensa puramente locale.

Due schemi di coordinamento ricorrono in letteratura
\citep{li2018drlsurvey}: l'\emph{addestramento centralizzato con
  esecuzione decentralizzata} (CTDE), in cui gli agenti condividono
informazione globale solo durante il training per poi decidere
localmente, e l'\emph{addestramento ed esecuzione entrambi
  decentralizzati} (DTDE), in cui l'informazione resta sempre e solo
locale, anche durante l'addestramento. Una scelta indipendente da questa
è se i diversi agenti condividano o meno gli stessi pesi (\emph{parameter
  sharing}): condividerli riduce il numero di parametri totali e permette a
un agente di beneficiare dell'esperienza raccolta da tutti gli altri, a
costo di una minore capacità di specializzarsi su condizioni locali
peculiari, a meno che il meccanismo non trovi un altro modo per
differenziare il comportamento tra agenti pur condividendo i pesi.

\textbf{Questo lavoro rappresenta un sistema multi-agente} in cui ogni intersezione della
rete stradale è un agente a sé, con un proprio stato $s_i^t$, una propria
azione $a_i^t$ (la fase semaforica scelta) e una propria ricompensa
$r_i^t$, calcolati indipendentemente per ogni intersezione $i$
(Capitolo~\ref{cap:tsc-modelli}). \textbf{L'architettura è di tipo DTDE}, con scambio di
informazione locale tra intersezioni vicine (§\ref{sec:background-gnn}), e
adotta il parameter sharing: tutte le intersezioni condividono la stessa
rete $Q$, ma un meccanismo di meta-learning genera pesi diversi per nodi con caratteristiche locali diverse, così da
attenuare proprio il costo della condivisione appena descritto.

\section{Reti Neurali Ricorrenti}
\label{sec:background-rnn}

Una rete neurale \emph{feedforward} calcola il proprio output come
funzione del solo input corrente. Dunque ogni esempio è trattato in modo
indipendente dai precedenti, e la rete non conserva alcuna memoria di
cosa abbia visto in passato. Questa assunzione fallisce ogni volta che i
dati hanno una struttura sequenziale in cui l'ordine e la storia contano
quanto il singolo valore osservato (testo, audio, serie temporali, o lo
stato di un'intersezione stradale). Le \emph{reti neurali ricorrenti} (RNN)
affrontano il problema introducendo un ciclo: a ogni istante $t$, un nodo
riceve in ingresso non solo il dato corrente $x^{t}$ ma anche il proprio
stato nascosto al passo precedente $h^{t-1}$,
\[h^{t} = \sigma\!\left(W_{hx} x^{t} + W_{hh} h^{t-1} + b_h\right),\label{eq:rnn-recurrence}\]
così che l'informazione rilevante di tutta la sequenza osservata fino a
$t$ possa, in linea di principio, persistere in $h^{t}$. Lo stesso
insieme di pesi $W_{hx}, W_{hh}, b_h$ viene riapplicato a ogni istante di
tempo, invece di avere pesi diversi per ogni passo. "Srotolando" la rete
lungo l'asse del tempo, cioè disegnando esplicitamente una copia dei nodi
per ogni istante $t=1,2,\dots$ collegata alla successiva (Figura~\ref{fig:rnn-unrolled}), ogni passo
diventa uno strato di una rete feedforward molto profonda a pesi
condivisi, addestrabile con la stessa backpropagation vista in precedenza
applicata lungo la sequenza anziché lungo i layer: questo algoritmo
prende il nome di \emph{backpropagation through time} (BPTT).

\begin{figure}[htbp]
  \centering
  \includegraphics[width=0.75\textwidth]{rnn_unrolled}
  \caption{La rete ricorrente srotolata lungo il tempo: il modulo $A$
    rappresenta la funzione ricorrente, ad esempio l'Equazione~\ref{eq:rnn-recurrence},
    riapplicata con gli stessi pesi a ogni istante.}
  \label{fig:rnn-unrolled}
\end{figure}

\paragraph{Il problema del gradiente.} Una rete ricorrente addestrata su
sequenze lunghe soffre di un problema strutturale: il gradiente
propagato a retroso attraverso molti passi temporali tende a annullarsi
(\emph{vanishing gradient}) o a divergere (\emph{exploding gradient}) a
seconda che i pesi ricorrenti abbiano modulo minore o maggiore di uno,
poiché il gradiente si ottiene moltiplicando ripetutamente per lo stesso
fattore a ogni passo srotolato, esattamente come una potenza $c^n$ tende a
zero o esplode al crescere di $n$ a seconda che $|c|$ sia minore o
maggiore di uno. Nella pratica, questo rende di fatto impossibile
apprendere dipendenze che coprono più di poche decine di passi con una
RNN semplice.

\paragraph{Long Short-Term Memory (LSTM).} L'LSTM \citep{lipton2015rnnreview}
risolve il problema sostituendo il semplice nodo ricorrente con una
\emph{cella di memoria}: uno stato interno $c^{t}$ collegato a se stesso
da un percorso a peso fisso unitario (il \emph{constant error carousel}),
attraverso cui il gradiente può fluire per molti passi senza annullarsi né
esplodere. L'accesso a questo stato è
regolato da tre \emph{gate}, unità sigmoidali il cui valore in $[0,1]$
modula quanta informazione passi attraverso di esse (Figura~\ref{fig:lstm-cell}).

\begin{figure}[htbp]
  \centering
  \includegraphics[width=0.7\textwidth]{lstm_cell}
  \caption{La cella di memoria dell'LSTM: il forget gate decide quanto
    dello stato precedente mantenere, l'input gate quanto del nuovo
    contenuto candidato incorporare, l'output gate quanto dello stato
    risultante esporre come output.}
  \label{fig:lstm-cell}
\end{figure}

A livello formale, con pesi $W_g \in \mathbb{R}^{2d_1 \times d_1}$ e bias $b_g \in
  \mathbb{R}^{d_1}$ condivisi da ogni nodo, per ciascuno dei quattro gate
$g \in \{\text{forget, input, output, cell}\}$, si ha:
\begin{equation}
  \begin{bmatrix} f_i^t \\ \iota_i^t \\ o_i^t \\ g_i^t \end{bmatrix} =
  \begin{bmatrix} \sigma \\ \sigma \\ \sigma \\ \tanh \end{bmatrix}
  \left(
  W_{\{f,\iota,o,g\}}^{\top} [\,e_i^t \,\|\, h_i^{t-1}\,] + b_{\{f,\iota,o,g\}}
  \right),
  \label{eq:lstm-background-gates}
\end{equation}
\[ c_i^t = f_i^t \odot c_i^{t-1} + \iota_i^t \odot g_i^t, \]
\[ x_i^t = h_i^t = o_i^t \odot \tanh(c_i^t); \]
il \emph{forget gate} $f$ decide quanto dello stato precedente
dimenticare, l'\emph{input gate} $\iota$ quanto del nuovo contenuto
candidato $g$ incorporare
nello stato aggiornato, e l'\emph{output gate} $o$ quanto dello stato
risultante esporre come output della cella. Un valore di forget gate
vicino a 1 lascia passare l'informazione pregressa quasi inalterata anche
dopo molti passi; uno vicino a 0 la cancella, permettendo alla rete di
"dimenticare" deliberatamente informazione non più rilevante.

\paragraph{Gated Recurrent Unit (GRU).} Una variante più leggera dell'LSTM,
la \textbf{GRU} \citep{cho2014gru}, fonde il \emph{forget gate} e
l'\emph{input gate} in un unico \emph{update gate} e rinuncia allo stato di
cella $c^{t}$ separato, esponendo un solo stato nascosto $h^{t}$ sia
come memoria interna sia come output. Si hanno così meno parametri da addestrare e un
grafo computazionale più semplice, a fronte di una capacità di
rappresentazione leggermente inferiore rispetto ai quattro gate distinti
dell'LSTM. Le due varianti raggiungono nella pratica prestazioni simili su
molti compiti, e la scelta tra loro è spesso guidata più dalla
convenzione di un'architettura di riferimento che da un vantaggio netto
dell'una sull'altra: questo lavoro adotta l'LSTM per coerenza con il
modello di riferimento \citep{wang2022metastgat}.

\section{Grafi e Graph Neural Network}
\label{sec:background-gnn}

\subsection{Cos'è un grafo}

Un \emph{grafo} $G = (V, E)$ è una struttura matematica composta da un
insieme di \emph{nodi} $V$ e un insieme di \emph{archi}
$E$, dove ogni arco collega una coppia di nodi. È lo strumento standard
per rappresentare entità e le relazioni tra loro, quando quelle relazioni
non hanno una struttura regolare o sequenziale. Ne sono esempi una rete di amicizie tra
persone, le pagine web collegate da link, le molecole i cui atomi sono
legati da legami chimici, e, in questo lavoro, una rete stradale, dove i
nodi sono le intersezioni e gli archi le strade che le collegano
direttamente.

Un arco può essere \emph{diretto} (se indica un verso) o \emph{non diretto} (se indica una relazione simmetrica, Figura~\ref{fig:graph-directed});
può inoltre avere un \emph{peso} associato, un numero che quantifica l'intensità o il costo della relazione.

\begin{figure}[htbp]
  \centering
  \includegraphics[width=0.75\textwidth]{graph_directed_undirected}
  \caption{Esempi di grafo diretto (sinistra) e non diretto (destra).}
  \label{fig:graph-directed}
\end{figure}

L'insieme dei nodi
collegati a un nodo $i$ tramite un arco si chiama
\emph{vicinato} di $i$, indicato $\mathcal{N}(i)$; il numero di archi
incidenti su $i$ è il suo \emph{grado}. Una sequenza di nodi collegati
uno al successivo da un arco è un \emph{cammino}; la lunghezza del
cammino più breve tra due nodi qualunque è la loro \emph{distanza} sul
grafo, misurata in numero di archi. Un grafo con $N$ nodi si può
rappresentare in forma matriciale con la \emph{matrice di adiacenza} $A
  \in \{0,1\}^{N \times N}$, dove $A_{iu}=1$ se e solo se esiste un arco tra
$i$ e $u$ (o il peso dell'arco, se il grafo è pesato, invece del semplice
$1$). È spesso utile
aggiungere a ogni nodo un \emph{self-loop}, un arco da un nodo verso se
stesso ($A_{ii}=1$), cosicché un nodo sia incluso nel proprio stesso
vicinato: la notazione $\hat A = A + I$ indica questa aggiunta e ricorre
nelle equazioni seguenti.

Un piccolo esempio: la matrice di adiacenza senza pesi del grafo non diretto in Figura~\ref{fig:graph-directed} è data da
\[A = \begin{pmatrix} 0&1&1 \\ 1&0&1 \\ 1&1&0 \end{pmatrix}.\]
Con l'aggiunta dei self-loop, $\hat A = A+I$ ha inoltre un $1$ su ogni elemento della diagonale.

Nella rete stradale di questo lavoro, ogni nodo $i$ del grafo porta con sé
un vettore di caratteristiche (o \emph{embedding}) che ne descrive lo
stato corrente, ad esempio quanti veicoli attendono su ciascuna corsia.
Il problema che le \emph{graph neural network} risolvono è come far sì che l'embedding di un nodo
tenga conto anche di ciò che accade nei nodi a lui collegati, non solo del proprio stato isolato.

\subsection{Message passing}

Le \emph{Graph Neural Network} (GNN) affrontano questo problema con un
principio comune, il \emph{message passing} \citep{wu2021gnnsurvey,
  bacciu2020gentlegnn}. A ogni iterazione (o \emph{layer}) $l$, ogni nodo
$i$ con rappresentazione corrente $h_i^{(l)}$ esegue due operazioni in
successione: prima un passo di \emph{aggregazione}, che combina i
messaggi ricevuti da ciascun vicino $u \in \mathcal{N}(i)$ in un unico
vettore di messaggio $m_i$,
\[m_i = \mathrm{AGGREGATE}\big(\{h_u^{(l)} : u \in \mathcal{N}(i)\}\big),\]
poi un passo di \emph{aggiornamento}, che combina il messaggio
aggregato con la rappresentazione precedente dello stesso nodo per
produrre quella nuova,
\[h_i^{(l+1)} = \mathrm{UPDATE}\big(h_i^{(l)},\ m_i\big).\]
Tutte le GNN discusse in questa tesi, incluse le due formulazioni di
seguito, sono istanze di questo stesso schema generale, che differiscono
solo nella scelta concreta di AGGREGATE e UPDATE. Nella GCN, AGGREGATE è
una media pesata dalla sola topologia; nella GAT, un peso appreso in base
al contenuto dei due nodi coinvolti. Impilando più layer, l'informazione
può viaggiare oltre il solo vicinato immediato: dopo $k$ layer, la
rappresentazione di un nodo dipende (indirettamente) da tutti i nodi
raggiungibili entro $k$ archi di distanza, poiché il messaggio ricevuto da
un vicino al layer $l$ incorpora già, ricorsivamente, i messaggi che
quel vicino aveva ricevuto dai propri stessi vicini al layer $l-1$.

\begin{figure}[htbp]
  \centering
  \includegraphics[width=0.85\textwidth]{gnn_pipeline}
  \caption{Lo schema generale di message passing impilato su più layer: a
    ogni layer ciascun nodo aggrega le rappresentazioni
    dei propri vicini e aggiorna la propria, prima di applicare la non
    linearità $\sigma$ e passare al layer successivo.}
  \label{fig:gnn-pipeline}
\end{figure}

\paragraph{Graph Convolutional Network (GCN).} \citet{kipf2017gcn}
propongono la formulazione più basilare di message passing tra quelle qui
discusse. La rappresentazione di ciascun nodo si aggiorna mediando quelle
dei vicini con un peso fisso, determinato unicamente dalla topologia
locale (quanti vicini ha un nodo), non dal contenuto delle
rappresentazioni stesse. Indicando con $\hat A = A + I$ la matrice di
adiacenza con self-loop e con $\hat D$ la matrice diagonale dei gradi
corrispondente ($\hat D_{ii} = \sum_u \hat A_{iu}$), la regola di
propagazione di un singolo layer si scrive, in forma matriciale:
\begin{equation}
  H^{(l+1)} = \sigma\!\left(\hat D^{-1/2} \hat A\, \hat D^{-1/2} H^{(l)} W^{(l)}\right)
  \label{eq:gcn-propagation}
\end{equation}
dove $H^{(l)} \in \mathbb{R}^{N \times d_l}$ raccoglie le rappresentazioni
di tutti gli $N$ nodi del grafo al layer $l$, $W^{(l)}$ è la matrice di
pesi condivisa da ogni nodo e $\sigma$ una non linearità puntuale. Per un
singolo nodo $i$, la normalizzazione simmetrica $\hat D^{-1/2}\hat
  A\hat D^{-1/2}$ equivale a pesare il contributo di ogni vicino $u$ con
$1/\sqrt{\hat d_i \hat d_u}$:
\begin{equation}
  h_i^{(l+1)} = \sigma\!\left(\sum_{u \in \mathcal{N}(i) \cup \{i\}}
  \frac{1}{\sqrt{\hat d_i \hat d_u}}\, W^{(l)} h_u^{(l)}\right)
  \label{eq:gcn-aggregate}
\end{equation}
un peso identico per ogni arco del vicinato di $i$, funzione solo del
grado dei due nodi coinvolti e mai del loro contenuto. È, per costruzione,
il termine di paragone più naturale rispetto a cui misurare cosa aggiunge
un meccanismo di attenzione appreso come quello discusso di seguito.

\paragraph{Graph Attention Network (GAT).} \citet{velickovic2018gat}
sostituiscono il peso fisso della GCN con un coefficiente di attenzione
appreso, specifico per ciascuna coppia di nodi collegati: il contributo di
un vicino non dipende più soltanto dalla topologia, ma anche dal suo
contenuto. Dato un nodo $i$ con vicinato $\mathcal{N}(i)$, embedding di
input $h_i$, matrice di proiezione condivisa $W$ e vettore di attenzione
$\vec a$:
\begin{equation}
  e_{iu} = \mathrm{LeakyReLU}\!\left(\vec a^\top [W h_i \,\|\, W h_u]\right),
  \qquad
  \alpha_{iu} = \frac{\exp(e_{iu})}{\sum_{u' \in \mathcal{N}(i)} \exp(e_{iu'})}
  \label{eq:gat-attention}
\end{equation}
\begin{equation}
  h_i' = \sigma\!\left(\sum_{u \in \mathcal{N}(i)} \alpha_{iu}\, W h_u\right)
  \label{eq:gat-aggregate}
\end{equation}
dove $\|$ denota la concatenazione: $e_{iu}$ è un punteggio di
compatibilità grezzo tra $i$ e $u$, che il \emph{softmax} sul vicinato
trasforma in un coefficiente $\alpha_{iu} \in [0,1]$ che somma a uno su
tutti i vicini di $i$, esattamente come una distribuzione di probabilità.

\begin{figure}[htbp]
  \centering
  \includegraphics[width=0.95\textwidth]{gat_attention}
  \caption{Il calcolo dell'attenzione per il nodo $1$ rispetto al proprio
    vicinato $\{1,2,3,4\}$ (con self-loop incluso): ogni coppia
    $(h_1,h_u)$ produce un coefficiente $\alpha_{1u}$, usato poi per
    combinare linearmente gli embedding dei vicini nella nuova
    rappresentazione $h_1'$. La figura, semplificata a scopo illustrativo,
    accorpa in un unico passo "Linear" il calcolo di compatibilità e la
    normalizzazione softmax descritti nel testo, e omette la proiezione
    condivisa $W$ e la non linearità $\sigma$.}
  \label{fig:gat-attention}
\end{figure}

Con $K$ teste di attenzione indipendenti (\emph{multi-head attention}),
ciascuna con i propri pesi $W$ e $\vec a$, gli output si concatenano nei
layer intermedi e si mediano nell'ultimo, secondo lo schema introdotto per
l'attenzione da \citet{vaswani2017attention} In questo modo teste diverse possono
imparare a specializzarsi su aspetti diversi della relazione tra due nodi.
A differenza della GCN, il coefficiente $\alpha_{iu}$ consente a un nodo
di ponderare diversamente i propri vicini in base al loro contenuto, non
solo alla loro posizione nel grafo: una proprietà rilevante per
un nodo i cui vicini abbiano importanza marcatamente diversa,
come nel caso di un incrocio stradale in cui passano un'arteria principale e una strada
secondaria. La funzione di compatibilità additiva appena descritta, per
quanto la più nota, non è l'unica possibile: il Capitolo~\ref{cap:tsc-modelli}
(Meta-GAT) ne impiega una variante strutturalmente diversa, con Query, Key
e Value esplicite e distinte al posto della proiezione $W$ condivisa, e un
prodotto scalare al posto della concatenazione.

\subsection{Invarianza e equivarianza per permutazione}
\label{subsec:permutation-invariance}

Sia l'aggregazione della GCN sia quella della GAT condividono una
proprietà che vale la pena rendere esplicita, perché è ciò che permette a
queste reti di funzionare su grafi di dimensione diversa senza dover
riaddestrare nulla. Entrambe calcolano il messaggio aggregato $m_i$ di un
nodo $i$ come una somma (pesata, uniformemente nella GCN o per attenzione
nella GAT) sui vicini $u \in \mathcal{N}(i)$: una somma non dipende
dall'ordine in cui i suoi addendi vengono elencati, quindi $m_i$ non
dipende dall'ordine in cui i vicini di $i$ vengono numerati o elencati, ma
solo da quali essi siano. Questa proprietà si chiama \emph{invarianza per
  permutazione} dell'operatore di aggregazione, ed è ciò che distingue
strutturalmente una GNN da una rete feedforward applicata a un vettore di
caratteristiche di lunghezza fissa: una MLP a cui si desse in input, in un
qualche ordine arbitrario, le caratteristiche di ciascun nodo del grafo
concatenate assegnerebbe implicitamente un significato specifico a ciascuna
posizione del vettore (il primo blocco di ingressi, il secondo, ...) e
produrrebbe quindi un output diverso se semplicemente si rinumerassero i
nodi, pur descrivendo esattamente la stessa rete stradale;
\citet{wei2019colight} osservano esattamente questo limite per gli schemi
di coordinamento basati su una lista ordinata di vicini
(§\ref{sec:background-tsc}). Applicando la stessa proprietà a ogni nodo del
grafo simultaneamente, l'intero layer di message passing risulta
\emph{equivariante per permutazione}: rinumerare i nodi del grafo in
ingresso produce in uscita esattamente le stesse rappresentazioni invece di un risultato arbitrariamente
diverso \citep{bronstein2021geometric}.

Questa proprietà ha una conseguenza pratica diretta per la generalizzazione
discussa nel Capitolo~\ref{cap:confronto} in quanto rende possibile la generalizzazione a topologie di reti stradali mai viste che il
capitolo stesso conduce su MetaSTGAT e sulle sue varianti.

\subsection{Reti spatio-temporali su grafo}

Molti grafi non sono statici: la topologia resta fissa, ma le
caratteristiche osservate a ciascun nodo evolvono nel tempo, e le due
dimensioni, spaziale e temporale, risultano intrecciate, poiché il valore
osservato a un nodo in un dato istante è correlato sia con i valori
passati allo stesso nodo sia con quelli ai nodi limitrofi. La famiglia di modelli che combina un layer
di propagazione sul grafo (GCN o GAT, appena visti) con una componente
ricorrente (LSTM, §\ref{sec:background-rnn}) sull'asse temporale prende il
nome di \emph{spatiotemporal graph neural network} (STGNN), un termine
ormai consolidato nella letteratura sul forecasting di serie temporali
multivariate correlate \citep{cini2025graphdl}. MetaSTGAT, descritto nel
dettaglio nel Capitolo~\ref{cap:tsc-modelli}, costituisce un esempio di
questa famiglia applicata al controllo semaforico anziché alla previsione:
un layer di \emph{graph attention} cattura la dipendenza spaziale tra
intersezioni limitrofe, un layer ricorrente (LSTM) ne cattura l'evoluzione
temporale, e i due vengono addestrati congiuntamente end-to-end.

\subsection{Pesi generati da meta-learning}

Le strutture neurali introdotte, GCN e GAT, nella loro formulazione originale, condividono gli
stessi pesi su tutti i nodi del grafo: un'assunzione ragionevole
quando i nodi sono sufficientemente omogenei, meno quando, come in una
rete stradale reale, la topologia locale di ciascun nodo (numero di
strade entranti, presenza di un'arteria, posizione periferica o centrale)
può differire sensibilmente da nodo a nodo. Un \emph{hypernetwork}
affronta il problema generando i pesi di un layer non come parametri
fissi appresi direttamente, bensì come output di una seconda rete,
condizionata sulle caratteristiche locali del nodo: è di fatto una via di
mezzo tra pesi completamente condivisi e pesi completamente indipendenti
per ogni nodo, che riduce il costo del parameter sharing senza rinunciare
del tutto ai suoi benefici. Il Capitolo~\ref{cap:tsc-modelli} descrive nel
dettaglio come MetaSTGAT applichi questo principio sia al layer di
\emph{graph attention} sia al layer ricorrente.

\subsection{Propagazione a lungo raggio}

Sia la GCN sia la GAT propagano informazione a un solo passo di distanza
(un arco nel grafo) per layer: raggiungere un vicino a distanza $k$
richiede impilarne altrettanti, con un costo di parametri crescente e un
rischio noto di \emph{over-smoothing}, per il quale le rappresentazioni
dei nodi divengono troppo omogenee al crescere della profondità (ogni nodo
finisce per "vedere" praticamente tutto il grafo, perdendo la capacità di
distinguersi dai propri vicini). Meccanismi più recenti, tra cui \textbf{SONAR} \citep{trenta2025sonar},
affrontano il problema della propagazione a lungo raggio secondo un
principio diverso, ispirato ai sistemi dinamici oscillatori piuttosto che
alla concatenazione sequenziale di layer. Il sistema si comporta come una
rete di masse collegate da molle e smorzatori, una massa per nodo del
grafo. Raccolto lo stato di tutti gli $N$ nodi al tempo continuo $t$ in
$X(t) \in \mathbb{R}^{N \times d}$, la sua evoluzione è governata da
un'equazione del second'ordine, un'accelerazione e non una semplice
velocità:
\begin{equation}
  \ddot X(t) = -\mathcal{L}^a X(t)\, W - D(X(t)) \odot \dot X(t) + F(X(t))
  \label{eq:sonar-wave-background}
\end{equation}
Tre termini, tre ruoli distinti. $\mathcal{L}^a$ è un Laplaciano di grafo
pesato da una conduttanza $a_{ui} \geq 0$ per arco, che tira lo stato di un
nodo verso quello dei suoi vicini $u \in \mathcal{N}(i)$, come una molla.
$D(\cdot) \geq 0$ è un termine dissipativo che agisce sulla velocità
$\dot X(t)$, frenando il movimento come un attrito. $F(\cdot)$ è una
forzante esterna, che inietta energia dall'esterno del sistema. Riscritta
come sistema del prim'ordine, con velocità ausiliaria $V=\dot X$, e
discretizzata con passo $h$ su $L$ iterazioni, l'aggiornamento a ogni
iterazione diventa:
\begin{equation}
  \begin{aligned}
    \big(\mathcal{L}^a X\big)_i & = \sum_{u \in \mathcal{N}(i)} a_{ui}\,(X_i - X_u),              \\
    V                           & \leftarrow V - h\big(\mathcal{L}^a X + D(X)\odot V - F(X)\big), \\
    X                           & \leftarrow X + h\,\tanh(V)
  \end{aligned}
  \label{eq:sonar-discrete-background}
\end{equation}
dove il $\tanh$ sull'aggiornamento di $X$ è un accorgimento di stabilità
numerica, assente nella derivazione originale ma presente
nell'implementazione ufficiale di \citet{trenta2025sonar}. Le $L$
iterazioni condividono sempre gli stessi pesi $a_{ui}$, $D(\cdot)$ e
$F(\cdot)$, per cui, a differenza di uno stack di $L$ layer indipendenti
come GCN o GAT, il costo di parametri di questo meccanismo resta costante
al crescere di $L$, un vantaggio strutturale dichiarato dagli stessi
autori. Nella formulazione originale, $a_{ui}$, $D(\cdot)$ e $F(\cdot)$
sono reti a pesi fissi, condivise da tutti i nodi, senza alcuna nozione di
meta-apprendimento.

\section{Traffic Signal Control come problema}
\label{sec:background-tsc}

\subsection{Tre generazioni di controllo}
\label{subsec:tsc-generations}

Il TSC si articola in due approcci, in ordine crescente di complessità e,
storicamente, di introduzione \citep{shams2018atsctaxonomy}.

\textbf{Controllo a tempi fissi.} \textsc{Fixed-Time} cicla le fasi
semaforiche di un incrocio in ordine fisso, cambiando fase a ogni step di
decisione, e non osserva mai lo stato reale della rete: non ha alcun
parametro da configurare oltre all'ordine del ciclo e non richiede
training. È estremamente semplice da implementare ed è tuttora la spina
dorsale della stragrande maggioranza dei sistemi semaforici reali, dove
le uniche modifiche diffuse riguardano l'ordine delle fasi (in cui una
fase può ripetersi più di una volta per ciclo), la durata di ciascuna
fase e un'estensione o restrizione del ciclo in base all'orario
giornaliero. Poiché non richiede sensori stradali in tempo reale né
potenza computazionale per decidere, è estremamente economico in termini
di strumentazione. I suoi svantaggi principali sono la rigidità e
l'incapacità di adattarsi alle fluttuazioni improvvise del traffico,
caratteristiche che possono risultare critiche in caso di eventi
imprevisti.

\textbf{Controllo reattivo.} Una regola locale osserva lo stato corrente
dell'intersezione (tipicamente il numero di veicoli in coda per corsia) e
sceglie l'azione che ottimizza un criterio istantaneo, senza affidarsi a
un piano prestabilito. Questa regola può essere scritta a mano, controllo
\emph{euristico}, oppure appresa dall'esperienza tramite reinforcement
learning, controllo \emph{appreso}: quest'ultimo è il paradigma adottato
da questo lavoro (§\ref{sec:background-rl}). All'interno del controllo
euristico, \citet{shams2018atsctaxonomy} distinguono ulteriormente il grado di
adattività: il controllo \emph{attuato} più semplice si limita a
prolungare o abbreviare le fasi di un piano di base ancora sostanzialmente
predefinito (estendere una fase finché un sensore rileva veicoli sulla
corsia servita, ad esempio), mentre un controllo pienamente
\emph{adattivo} ricalcola la struttura stessa del piano, non solo la sua
durata, in risposta al traffico osservato. Piattaforme come SCATS e SCOOT
si collocano in questa seconda fascia, coordinando intere aree urbane
aggiustando periodicamente i piani semaforici sulla base del traffico
osservato, pur restando ancorate a una struttura di piano pre-esistente
\citep{wei2020tscsurvey}.

\textsc{Max-Pressure} \citep{varaiya2013maxpressure} rappresenta invece
un sistema di controllo reattivo che garantisce matematicamente il
massimo throughput (la capacità di smaltimento dei veicoli) della rete
stradale senza alcuna necessità di training, partendo dalla teorica
assunzione che le corsie abbiano capacità infinita. Questo risultato, pur
non potendo essere traslato integralmente in scenari reali dove non è
rara la possibilità di una completa saturazione di una via, si fonda
solamente su dati raccolti a ogni istante decisionale nelle corsie
adiacenti all'incrocio e su un principio di base: massimizzare la
pressione di ciascuna fase semaforica. Data una fase semaforica $\phi$,
la sua pressione è definita come la somma, su tutti i movimenti (coppia
corsia di ingresso/corsia di uscita) abilitati da $\phi$, della
differenza tra il numero di veicoli in coda sulla corsia di ingresso e
quelli sulla corsia di uscita, contati sull'intera corsia:
\begin{equation}
  P(\phi) = \sum_{(\ell_{\mathrm{in}}, \ell_{\mathrm{out}}) \,\in\, \mathrm{movimenti}(\phi)} \big(n(\ell_{\mathrm{in}}) - n(\ell_{\mathrm{out}})\big),
  \label{eq:maxpressure-phase}
\end{equation}
dove $n(\ell)$ indica il numero di veicoli sulla corsia $\ell$. A ogni
step di decisione, l'algoritmo calcola la pressione di tutte le fasi e
seleziona quella a pressione massima, $\phi^* =
  \operatorname*{argmax}_\phi P(\phi)$, senza alcun piano né ciclo
predefinito, né alcun vincolo sulle azioni consecutive: la stessa fase
può in linea di principio essere selezionata un numero indefinito di
volte. Uno svantaggio è la necessità di sensoristica capace di monitorare
le corsie per la loro interezza. Il problema principale, tuttavia, è
l'assenza di garanzie di equità: in presenza di un forte afflusso di
traffico su un certo insieme di corsie, la pressione maggiore può
trovarsi sempre sulle fasi che coinvolgono quelle vie, rendendo
impossibile alle altre fasi la possibilità di essere selezionate, con
conseguenti tempi di attesa eccessivi per i veicoli sulle vie meno
trafficate, un aspetto di equità che questo lavoro esamina esplicitamente.

\textbf{Fixed-Time e Max-Pressure sono le due baseline sperimentali
  classiche a cui questo lavoro confronta i propri modelli}
(Capitolo~\ref{cap:confronto}), usate qui esattamente nella forma appena
descritta. Pur fondandosi su paradigmi opposti, statico l'uno e adattivo
l'altro, entrambe non richiedono alcuna fase di addestramento, il che le
rende candidati naturali per il confronto con i modelli di deep
reinforcement learning discussi nel resto di questo lavoro.

\paragraph{Equità nella letteratura recente.} Alcuni metodi affrontano
esplicitamente il problema dell'equità penalizzando le fasi responsabili
di attese eccessive. FELight, per esempio, introduce una metrica di
equità basata sui tempi di attesa dei veicoli e penalizza le fasi già
identificate come responsabili di un ritardo superiore a una soglia
prestabilita \citep{xiao2025tscreview}.

\paragraph{Generalizzazione nella letteratura recente.} La prassi più
diffusa nel controllo semaforico appreso è addestrare e valutare un
modello sulla stessa rete stradale, al più ripetendo l'intero
esperimento, addestramento incluso, su una rete diversa, non trasferendo
un singolo modello già addestrato a una rete che non ha mai visto
\citep{wei2020tscsurvey}. CoLight \citep{wei2019colight}, per esempio,
conduce cinque esperimenti indipendenti, due su reti sintetiche e tre su
reti reali, ciascuno con il proprio addestramento; MetaSTGAT
\citep{wang2022metastgat} segue lo stesso schema su una griglia sintetica
e due dataset reali. Quando in questi lavori compare il termine
\emph{generalization}, indica quasi sempre la tenuta a diverse intensità
di traffico sulla stessa rete, non a una topologia strutturalmente
diversa \citep{wang2026astgcndrl, sun2026tgmaddpg}, come reso esplicito
da Zhang et al. nella motivazione di GeneraLight
\citep{zhang2020generalight}. Le eccezioni restano circoscritte a
un'unica dimensione di variazione: AttendLight
\citep{oroojlooy2020attendlight} addestra su un insieme eterogeneo di
intersezioni singole e ne verifica le prestazioni su intersezioni mai
viste, ma dichiara che l'estensione a una rete di più intersezioni
coordinate resta lavoro futuro; Advanced-XLight
\citep{zhang2022advancedxlight} misura un rapporto di trasferimento tra
un modello addestrato su un dataset reale e valutato direttamente su un
altro, senza ulteriore addestramento, ma sempre tra reti di scala
paragonabile; lo stesso GeneraLight generalizza al variare del flusso di
traffico, non della topologia della rete.

\paragraph{Il meccanismo di meta-learning di MetaSTGAT.} MetaSTGAT
\citep{wang2022metastgat} combina la GNN e la componente ricorrente con
un modulo di \emph{meta-learning}: invece di condividere gli stessi pesi
tra tutte le intersezioni della rete, questo modulo genera dinamicamente
i pesi di uno strato di attenzione e di uno strato ricorrente a partire
da caratteristiche locali dell'intersezione, così che nodi diversi
possano essere serviti da trasformazioni diverse pur restando un unico
modello addestrato end-to-end. Gli autori originali riportano una
riduzione del travel time medio rispetto a CoLight \citep{wei2019colight},
un metodo di riferimento tra i più citati per il Traffic Signal Control e
basato su \emph{graph attention} ma con pesi condivisi.

\subsection{Simulazione}

Valutare una politica di controllo semaforico su una rete reale non è
praticabile durante lo sviluppo, in quanto cambiare i tempi di un semaforo cittadino
migliaia di volte per tentativi, come richiederebbe l'addestramento di un
agente RL, non è ovviamente proponibile. Serve quindi un simulatore che
riproduca la dinamica del traffico con un dettaglio sufficiente, pur
mantenendo un costo computazionale accettabile per l'addestramento di
centinaia di episodi.

Un simulatore di traffico microscopico rappresenta ogni veicolo
individualmente, con una propria posizione, velocità e traiettoria,
aggiornate a ogni istante da un modello che regola la velocità di un
veicolo in base al comportamento di quello che lo precede mantenendo una distanza di sicurezza.
È il livello di dettaglio presente in questo lavoro, sufficientemente fedele da catturare l'effetto di una singola
decisione semaforica su una coda specifica.

\textsc{CityFlow} \citep{zhang2019cityflow} è il simulatore microscopico scelto in questo lavoro. Rispetto ad alternative più diffuse
come SUMO \citep{lopez2018sumo}, \textsc{CityFlow} semplifica ulteriormente il modello che regola la dinamica dei veicoli per privilegiare la velocità di simulazione,
un compromesso adatto quando l'obiettivo primario è addestrare un agente su un numero elevato di episodi piuttosto
che riprodurre con precisione millimetrica la dinamica di un incrocio specifico.